\section{\label{sec:defs}分布函数}


本文通篇聚焦于味道非单态的无极化准 PDFs 与光前 PDFs。向极化准 PDFs 的推广是直接的。味道单态情形会引入与胶子分布的额外混合，但原理类似。我们还假定夸克是无质量的，并忽略正确处理重味所带来的复杂之处——对光前 PDFs 而言这仍是持续研究的课题（综述可例如参见
\cite{Thorne:2008xf,Olness:2008px,Binoth:2010nha}）。

\subsection{裸 PDFs}

在下文中，凡使用“裸”一词，均指在有限截断下由某种调节子确定的有限矩阵元。本节的讨论中我们把调节子隐含不写，不过读者心中不妨以维度正规化为例。这些裸矩阵元需要先在某个方案中重正化，然后才能移除调节子（或在格点上取连续极限）。这一用法沿用文献
\cite{Collins:1981uw,Collins:2011zzd} 中对光前 PDFs 的广泛讨论。

我们记裸光前 PDF 为
$f^{(0)}(\xi)$。光前 PDFs 常写作
$f_{j/N}^{(0)}(\xi)$，其中 $j$
标记夸克味、$N$ 标记核子种类；此处我们只考虑非单态分布，故可忽略部分子种类间的混合，并且我们的讨论足够普适，核子种类与本讨论无关。我们采用光前坐标
$(x^+,x^-,\boldsymbol{x}_\mathrm{T})$，满足 $x^\pm = (t\pm z)/\sqrt{2}$，
并定义 $\xi=k^+/P^+$。使用 $\xi$ 而不用 Bjorken-$x$，是为了与表征散射实验运动学的 Bjorken-$x$ 参数区分开；后者由实验动量转移 $Q^2=-q^2$ 和强子动量 $P$
表示为 $x=Q^2/(2P\cdot q)$。裸 PDF 定义为 \cite{Collins:2011zzd}
\begin{equation}
f^{(0)}(\xi) = \int_{-\infty}^\infty
\frac{\mathrm{d}\omega^-}{4\pi}
e^{-i\xi P^+\omega^-} \left\langle P \left|
T\,\overline{\psi}(0,\omega^-,\mathbf{0}_\mathrm{T})
W(\omega^-,0)\gamma^+\frac{\lambda^a}{2}\psi(0) \right|
P\right\rangle_\mathrm{C}.
\end{equation}
这里 $T$ 是
时序算符，$\psi$ 是夸克
场，下标 C 表示真空
期望值已作减除（换言之，只包含连通
贡献）。算符 $W(\omega^-,0)$ 是 Wilson 线，
\begin{equation}
W(\omega^-,0) =
{\cal P}\exp\left[-ig_0\int_0^{\omega^-}\mathrm{d}y^-A^+_\alpha(0,y^-,
\mathbf{0}_{\mathrm{T}})T_\alpha\right],
\end{equation}
其中 ${\cal P}$ 是路径排序算符，$g_0$ 是 QCD 裸耦合，
而 $A^\mu = A^\mu_\alpha T_\alpha$ 是带生成元
$T_\alpha$ 的 $SU(3)$ 规范势（对色指标 $\alpha$ 求和隐含）。靶态
$|P\rangle$ 是自旋平均的、具有相对论归一化的精确动量本征态：
\begin{equation}
\langle P' | P \rangle =
(2\pi)^32P^+\delta\left(P^+-P^{\prime\,+}\right)\delta^{(2)}
\left(\mathbf{P}_\mathrm{T} - \mathbf{P}'_\mathrm{T}\right).
\end{equation}

裸 PDFs 的矩定义为
\begin{equation}
a_{0}^{(n)} = \int_0^1
\mathrm{d}\xi\,\xi^{n-1}\left[
f^{(0)}(\xi)+(-1)^n\overline{f}^{(0)}(\xi)\right] = \int_{-1}^1
\mathrm{d}\xi\,\xi^{n-1}
f(\xi),
\end{equation}
其中 $\overline{f}^{(0)}(\xi)$ 是反夸克 PDF，第二个等号
源于夸克与反夸克 PDFs 之间的关系
\begin{equation}\label{eq:pdf2apdf}
f^{(0)}(-\xi) = -\overline{f}^{(0)}(\xi),
\end{equation}
只要夸克与反夸克场是经典的，或采用光前量子化，该式对裸分布便成立 \cite{Collins:1981uw}。

我们可以把这些裸矩 $a_{0}^{(n)}$ 经由
\begin{equation}\label{eq:opePDF}
\left\langle P | {\cal O}_{0}^{\{\mu_1\ldots \mu_n\}} | P
\right\rangle = 2a_{0}^{(n)} \left(P^{\mu_1}\cdots P^{\mu_n} -
\mathrm{traces}\right).
\end{equation}
与扭度2算符的矩阵元联系起来。这里的裸
扭度2算符为
\begin{equation}\label{eq:baretwist2def}
{\cal
O}_0^{\{\mu_1\cdots \mu_n\}} = i^{n-1}\overline{\psi}(0)
\gamma^{\{\mu_1}D^{\mu_2}\cdots
D^{\mu_n\}}\frac{\lambda^a}{2}\psi(0)-\operatorname{traces}.
\end{equation}
在这些表达式中，花括号表示对称化，$D^\mu$ 是
对称协变
导数，$\lambda^a$ 是 $SU(2)$ 味矩阵，减去迹项保证算符
在 $SU(2)_{\mathrm{L}}\otimes SU(2)_{\mathrm{R}}$ 下不可约地变换。

\subsection{重正化的 PDFs}

至此我们考虑的都是裸光前 PDFs，默认这类对象是用某种使
裸分布有限的调节子求值的。现在引入
重正化光前 PDFs。我们强调，本节考虑的是
保持转动对称性的重正化方案；为明确起见，
可以设想为 $\ms$ 方案。若使用破坏转动不变性的调节子（如
格点调节子），则会出现额外的复杂性。
由于我们将避免在有限格距下显式计算矩，这里不讨论此类复杂性。只要
不存在幂发散混合，所有在格点上计算的关联函数都可以
重正化并外推到连续极限。下一节中，我们将提出一种
不具有幂发散混合的涂抹
关联函数。

一般地，重正化光前 PDFs 借助核 ${\cal Z}(\zeta/\xi,\mu)$ 写作
\begin{equation}
f(\xi,\mu) = \int_\xi^1 \frac{\mathrm{d}\zeta}{\zeta}{\cal
Z}\left(\frac{\zeta}{\xi},\mu\right)f^{(0)}(\zeta),
\end{equation}
其中 $\mu$ 是某个重正化
能标。对于非单态分布，无需考虑部分子种类间的混合。用重正化光前 PDF 表出的
重正化 Mellin 矩为
\begin{equation}
a^{(n)}(\mu) = \int_0^1 \mathrm{d}\xi\,\xi^{n-1}\left[
f(\xi,\mu)+(-1)^n\overline{f}(\xi,\mu)\right] = \int_{-1}^1
\mathrm{d}\xi\,\xi^{n-1}
f(\xi,\mu),
\end{equation}
它可以经由
\begin{equation}\label{eq:RopePDF}
\left\langle P | {\cal O}^{\{\nu_1\ldots \nu_n\}}(\mu) | P
\right\rangle = 2a^{(n)}(\mu) \left(P^{\nu_1}\cdots P^{\nu_n} -
\mathrm{traces}\right).
\end{equation}
与重正化扭度2算符
${\cal O}^{\{\nu_1\ldots \nu_n\}}(\mu) =Z_{\cal O}(\mu){\cal
O}_0^{\{\nu_1\ldots \nu_n\}}$ 的矩阵元联系起来。只要光前 PDFs 与扭度2算符在同一方案中重正化，此关系即成立 \cite{Collins:1981uw}。



\subsection{涂抹准 PDFs}

我们通过对费米子场和规范场同时施加梯度流涂抹，构造有限的准 PDF 矩阵元
\cite{Narayanan:2006rf,Luscher:2011bx,Luscher:2013cpa}。梯度流是原始夸克与胶子场在新维度（通常称为“流时间”）中朝向 QCD 作用量经典极小的确定性演化 \cite{Luscher:2011bx,Luscher:2013cpa}。流时间演化的选取旨在消除紫外涨落，这相当于在实空间中将夸克与胶子场涂抹开来，涂抹尺度正比于流时间的平方根。此处我们不再详细描述梯度流，
读者可参阅近期综述
\cite{Luscher:2013vga,Ramos:2015dla} 以了解细节与应用。

就我们的目的而言，只需用到梯度流的下列性质。第一，梯度流充当规范不变的紫外
调节子。第二，给定零流时间处的重正化理论，涂抹
场的矩阵元自动有限，至多相差
费米子场的乘法性
波函数重正化 \cite{Makino:2014taa,Hieda:2016lly}。
第三，只要流时间以物理单位固定，涂抹场的格点矩阵元
在连续极限下保持有限 \cite{Luscher:2013cpa,Monahan:2015lha}。本质上，梯度流
使人可以用一个新的涂抹尺度调节子取代格点调节子。这最后一个事实使人能够确定诸如扭度2算符等
格点矩阵元
的连续极限，而不出现幂发散混合。在连续时空中，由于梯度流保持转动对称性，扭度2算符之间的混合退化为普通混合，系数
只依赖涂抹尺度而非格距倒数之幂
\cite{Monahan:2015lha}。

我们把流时间 $\tau$ 处的带环（ringed）费米子场记为
$\overline{\chi}(x;\tau)$ 与 $\chi(x;\tau)$，并把同一流时间处由涂抹规范场
$B_\mu(x;\tau)$ 构造的 Wilson 线记为 ${\bf\cal W}(0,z;\tau)$。
我们从矩阵元
\begin{equation}\label{eq:hdef}
h^{(s)}\left(\frac{z}{\sqrt{\tau}},\sqrt{\tau}P_z,
\sqrt{\tau}\Lambda_{\mathrm{QCD}},\sqrt{\tau}M_\mathrm{N}\right) =
\frac{1}{2P_z}\left\langle
P_z \left| \overline{\chi}(z;\tau)
{\bf\cal W}(0,z;\tau)\gamma_z\frac{\lambda^a}{2}\chi(0;\tau)\right|
P_z\right\rangle_\mathrm{C},
\end{equation}
出发；它是无量纲的，因而只依赖各尺度的无量纲组合。我们注意到流时间具有长度平方的量纲。下标 C
表示该矩阵元的非连通贡献已被去除。带环费米子场不需要波
函数重正化，因此只要流时间 $\tau$ 非零且以物理单位固定，这个涂抹矩阵元就是有限的，
因为由涂抹场构造的关联函数是有限的
\cite{Luscher:2011bx,Luscher:2013cpa}。注意当流时间趋于零时会出现发散，
届时矩阵元将需要
重正化。

随后我们把准 PDF
\cite{Ji:2013dva,Ji:2014gla} 定义为
\begin{equation}\label{eq:qpdfdef}
q^{\,(s)}\left(\xi,\sqrt{\tau}P_z,\sqrt{\tau}\Lambda_{\mathrm{QCD}},
\sqrt{\tau}M_\mathrm{N}\right)
=
\int_{-\infty}^\infty \frac{\mathrm{d}z}{2\pi} e^{i\xi z
P_z} P_z\,
h^{(s)}(\sqrt{\tau}z,\sqrt{\tau}P_z,\sqrt{\tau}\Lambda_{\mathrm{QCD}},
\sqrt{\tau}M_\mathrm{N}),
\end{equation}
其中 $\xi$ 是无量纲参数，可以被朴素地解释为
部分子在核子 $N$ 中携带的纵向动量分数。然而这一诠释在欧几里得空间并不正确；应把 $\xi$
视为傅里叶变换中的一个无量纲动量变量。

实践中，涂抹矩阵元 $h$ 是在有限格距 $a$ 的格点计算中确定的，即
\begin{equation}\label{eq:hacont}
h^{(s)}\left(\frac{z}{\sqrt{\tau}},\sqrt{\tau}P_z,\sqrt{\tau}\Lambda_{\mathrm{QCD}},
\sqrt{\tau}
M_\mathrm{N}\right) =  \lim_{a\to 0}
h\left(\frac{z}{a},\frac{\sqrt{\tau}}{a},aP_z,a\Lambda_{\mathrm{QCD}},aM_\mathrm{N}
\right) ,
\end{equation}
其中 $\sqrt{\tau}$ 与 $P_z$ 保持固定，而
\begin{align}
h\Big(\frac{z}{a},\frac{\sqrt{\tau}}{a},aP_z,a\Lambda_{\mathrm{QCD}},{} &
aM_\mathrm{N}\Big) = \nonumber\\
{} & \frac{1}{2aP_z}
\left\langle
aP_z \left|
\overline{\chi}\left(\frac{z}{a};\frac{\sqrt{\tau}}{a}\right)
W\left(0,\frac{z}{a};\frac{\sqrt{\tau}}{a}\right)\gamma_z\frac{\lambda^a}{2}
\chi\left(0;\frac{\sqrt{\tau}}{a}\right)\right|
aP_z\right\rangle_\mathrm{C}.
\end{align}
